\section{Residuation}
The key step is to show that product and division satisfy the expected adjoint behaviour.

\begin{proposition}[Right residuation]
For all formulae $A,B,C$,
\[
A\concat B \Sem C
\quad\Longleftrightarrow\quad
A \Sem C/B.
\]
\end{proposition}

\begin{proof}
Suppose first that $A\concat B \Sem C$. Let $x \in \Frag(\mathcal C)$. We must show that
$v(x,A) \leq v(x,C/B)$.

If the defining set for $v(x,C/B)$ contributes no witness beyond the default element $1$, then
$v(x,C/B)=1$ and the claim is immediate because $v(x,A) \in [0,1]$.

Otherwise take any $y \in \Frag(\mathcal C)$ such that $xy=w \in \Frag(\mathcal C)$ and
$v(y,B)>0$. Since $xy=w$, the decomposition $w=xy$ is one of the candidates in the product
clause, hence
\[
v(w,A\concat B) \geq v(x,A)\,v(y,B).
\]
By the assumption $A\concat B \Sem C$, we have $v(w,A\concat B) \leq v(w,C)$. Therefore
\[
v(x,A)\,v(y,B) \leq v(w,C).
\]
Because $v(y,B)>0$, division yields
\[
v(x,A) \leq \frac{v(w,C)}{v(y,B)}.
\]
As also $v(x,A)\leq 1$, we obtain
\[
v(x,A) \leq \min\Bigl(1,\frac{v(w,C)}{v(y,B)}\Bigr).
\]
This holds for every admissible witness $y$, so after taking the minimum in the definition of
$v(x,C/B)$ we conclude that $v(x,A) \leq v(x,C/B)$.

Conversely, assume $A \Sem C/B$. Let $w \in \Frag(\mathcal C)$. We show that
$v(w,A\concat B) \leq v(w,C)$. Consider any decomposition $w=xy$ with $x,y \in \Frag(\mathcal C)$.
If $v(y,B)=0$, then
\[
v(x,A)\,v(y,B)=0 \leq v(w,C).
\]
If $v(y,B)>0$, then by the assumption $A \Sem C/B$ we have
\[
v(x,A) \leq v(x,C/B).
\]
Now the definition of $v(x,C/B)$ is a minimum over admissible witnesses, one of which is
precisely the present decomposition $xy=w$. Hence
\[
v(x,C/B) \leq \min\Bigl(1,\frac{v(w,C)}{v(y,B)}\Bigr).
\]
Combining the two inequalities and multiplying by $v(y,B)$ gives
\[
v(x,A)\,v(y,B) \leq v(w,C).
\]
Thus every candidate product score for a split of $w$ is bounded above by $v(w,C)$, and so
its maximum is as well:
\[
v(w,A\concat B) \leq v(w,C).
\]
Therefore $A\concat B \Sem C$.
\end{proof}

\begin{proposition}[Left residuation]
For all formulae $A,B,C$,
\[
A\concat B \Sem C
\quad\Longleftrightarrow\quad
B \Sem A\backslash C.
\]
\end{proposition}

\begin{proof}
The proof is symmetric to the previous one. Instead of decompositions of the form $xy=w$, one
uses decompositions of the form $yx=w$, which are exactly the contexts quantified over in the
definition of $v(x,A\backslash C)$.
\end{proof}

